\documentclass[10pt,a4paper]{amsart}
\usepackage[margin=19mm]{geometry}
\usepackage{amsmath,amssymb,mathtools}
\usepackage[scr=rsfs,cal=euler]{mathalfa}
\usepackage{xfrac}
\usepackage[unicode,hidelinks,bookmarks=false]{hyperref}
\newcommand{\ie}{i.e.\ }
\newcommand{\cf}{cf.\ }
\newcommand{\resp}{respectively\ }
\newcommand{\Subsection}{\subsection}
\providecommand{\nobreakdash}{\nobreak\hbox{-}}
\setlength{\emergencystretch}{2em}
\multlinegap0pt
\usepackage{amssymb}
\usepackage{enumerate}
\usepackage{xcolor}
\usepackage{tikz}
\newtheorem{lemma}{Lemma}
\newcommand{\gos}{GOS}
\title{Integral stochastic orders of $m$-generalized order statistics from transform-ordered nonparametric families}
\author{Idir Arab, Tommaso Lando, Paulo Eduardo Oliveira, Tomasz Rychlik}
\date{Source: arXiv:2606.07022v1 (CC BY 4.0)}
\begin{document}
\maketitle

\noindent\emph{Source and licence.} Integral stochastic orders of $m$-generalized order statistics from transform-ordered nonparametric families. Version arXiv:2606.07022v1 (CC BY 4.0), distributed by arXiv under the Creative Commons Attribution 4.0 International licence, http://creativecommons.org/licenses/by/4.0/, as recorded in the arXiv licence metadata for this exact version and shown on the abstract page https://arxiv.org/abs/2606.07022v1. Full licence notice: \texttt{licenses/arxiv-2606.07022-CC-BY-4.0.txt}.\par\smallskip

\noindent\textit{Editorial scope.} Counted unit: the article's lemma t11 (source0.tex lines 604--608). The article's complete original proof of it is delivered with it (source0.tex lines 609--685, one \texttt{proof} environment of 77 lines). No mathematics is written, repaired or re-derived: the statement and the proof are the article's own lines, in the article's own notation, and no retained line is substituted or re-flowed.  No further source-local context is needed: every cross-reference of every retained line resolves inside the two delivered spans.  Nothing was omitted from any retained span, so the package carries no omission record.  No retained line contains a citation, so the wrapper carries no bibliography and no bibliography entry.  No retained line contains a figure environment, an image-inclusion command or drawing code, so no image is delivered and the package declares no asset dependency.  Editorial formatting only, in addition: an AMS \texttt{amsart} preamble that restates the article's own declarations and macros used by the retained lines, namely the environments \texttt{lemma} on separate counters, together with the standard packages those lines need.  The retained environments print in this document as Lemma 1 in order of appearance, whereas the article prints the same items with the numbering of its own full document.  The complete native source of the article as distributed in the arXiv e-print for this exact version is bundled unchanged as \texttt{sources/202227/source0.tex}.
\begin{lemma}\label{t11}
Let $f_{r,\gamma_{1:r},\mu}$ and $f_{q,\beta_{1:q},\mu}$ denote the density functions of the $r$-th  and $q$-th, $r \neq q$, uniform $m$-\gos, respectively,  
with the minimal parameters $\gamma_{1:r} > \beta_{1:r}>0$, respectively, and identical common differences $\mu \geq 0$. 
If $r>q$ then the difference $f_{r,\gamma_{1:r}, \mu} - f_{q, \beta_{1:q}, \mu}$ has the sign variation $-+-$ in $(0,1)$. If $r<q$, then this difference has sign variation $+-$ in $(0,1)$.
\end{lemma} 

\begin{proof}
Suppose that $\mu>0$ and $r>q$. Then
\begin{eqnarray}
\label{14a}
\lefteqn{f_{r, \gamma_{1:r}, \mu} (x) - f_{q, \beta_{1:q}, \mu} (x)
=  f_{q, \beta_{1:r}, \mu} (x)} \nonumber \\
& \qquad\times & \left[
\frac{M(r,\gamma_{1:r},\mu)}{M(q,\beta_{1:q},\nu)}
\frac{(q-1)!}{(r-1)!}\frac{1}{\mu^{r-q}}
(1-x)^{\gamma_{1:r}-\beta_{1:r}} \bigl[1- (1-x)^\mu\bigr]^{r-q}-1\right]\!.
\end{eqnarray}
Under change of variables $y= (1-x)^\mu$ the expression in the second line has the form\linebreak 
$c y^{\frac{\gamma_{1:r}-\beta_{1:r}}{\mu}}(1-y)^{r-q}-1$ for a positive $c$. It takes the value $-1$ when $y=0$ or $y=1$, and it is increasing-decreasing in between, because its derivative
\[
c y^{\frac{\gamma_{1:r}-\beta_{1:r}}{\mu}-1}
(1-y)^{r-q-1}\left[ \frac{\gamma_{1:r}-\beta_{1:r}}{\mu} (1-y) -
(r-q)y \right]
\]
is first positive and then negative.
Consequently, (\ref{14a}) is either negative, positive and negative in $(0,1)$ or merely negative there. However, the latter case is impossible.

\smallskip

Consider now the case when $\mu=0$ and $r>q$. We now have
%\begin{eqnarray}
%\lefteqn{f_{r, \gamma_{1:r}, 0} (x) - f_{q, \beta_{1:r}, 0} (x)} \nonumber \\
% & = & f_{q, \beta_{1:r}, 0} (x)   \left[ \frac{\gamma_{1:r}^r}{\beta_{1:r}^q}
%\frac{(q\!-\!1)!}{(r\!-\!1)!} (1-x)^{\gamma_{1:r}-\beta_{1:r}}
%(-\ln (1-x))^{r-q}-1 \right].
%\end{eqnarray}
\begin{equation}
\label{14b}
f_{r, \gamma_{1:r}, 0} (x) - f_{q, \beta_{1:r}, 0} (x)
  =  f_{q, \beta_{1:r}, 0} (x)   \left[ \frac{\gamma_{1:r}^r}{\beta_{1:r}^q}
\frac{(q\!-\!1)!}{(r\!-\!1)!} (1-x)^{\gamma_{1:r}-\beta_{1:r}}
(-\ln (1-x))^{r-q}-1 \right].
\end{equation}
The expression in the brackets is equal to $-1$ at $0$ and $1$, and it is increasing-decreasing in $(0,1)$, because its
derivative
\[
\frac{\gamma_{1:r}^r}{\beta_{1:r}^q}
\frac{(q-1)!}{(r-1)!} (1-x)^{\gamma_{1:r}-\beta_{1:r}-1}
\Bigl[-\ln (1-x)]^{r-q-1}[(\gamma_{1:r}-\beta_{1:r})\ln (1-x) +r-q\Bigr]
\]
is first positive and then negative.  In conclusion, the sign variation of (\ref{14b}) is $-+-$.

\smallskip

We proceed now to the case $r<q$. For $\mu>0$ yields
\begin{eqnarray}
\lefteqn{f_{r, \gamma_{1:r}, \mu} (x) - f_{q, \beta_{1:q}, \mu} (x)
=  \frac{M(r,\beta_{1:r},\nu)}
{(q-1)!\mu^{r-1}}
(1-x)^{\beta_{1:r}-1} [1-(1-x))^\mu ]^{r-1}} \nonumber \\
& & \qquad \times \left[ 
 \frac{M(r,\gamma_{1:r},\mu)}{M(q,\beta_{1:q},\nu)}
\frac{(q-1)!}{(r-1)!}
(1-x)^{\gamma_{1:r}-\beta_{1:q}} - \frac{1}{\mu^{q-r}}
[1-(1-x)^\mu ]^{q-r} \right].
\label{14c}
\end{eqnarray}
The first term in the brackets decreases from a positive number when $x=0$ to $0$ when $x=1$, whereas the following subtrahend increases from $0$ when $x=0$ to a positive value when $x=1$. Therefore (\ref{14c}) is $+-$ on $(0,1)$, as claimed.

We finally consider the case $\mu=0$ with $r<q$. We have
\begin{eqnarray}
\lefteqn{f_{r, \gamma_{1:r}, 0} (x) - f_{q, \beta_{1:q},0} (x)
=  \frac{\beta_{1:r}^q}{(q-1)!}
(1-x)^{\beta_{1:q}-1} [-\ln(1-x)]^{r-1}} \nonumber \\
& & \qquad\qquad\quad \times \left[ \frac{\gamma_{1:r}^r}{\beta_{1:q}^q}
\frac{(q-1)!}{(r-1)!}
(1-x)^{\gamma_{1:r}-\beta_{1:q}} -
\bigl[-\ln (1-x) \bigr]^{q-r} \right].
\label{14d}
\end{eqnarray}
The first term in the second line decreases from a positive number at $0$ to $0$ at $1$, and the latter one increases from $0$ to $+\infty$
on $(0,1)$. Therefore the sign variation of (\ref{14d}) is $+-$. This completes the proof of Lemma \ref{t11}.
\end{proof}

\end{document}
